% Part: normal-modal-logic
% Chapter: axioms-systems
% Section: consistency

\documentclass[../../../include/open-logic-section]{subfiles}

\begin{document}

\olfileid{nml}{prf}{con}

\olsection{Konsistensi}

Konsistensi iku sipat wigati saka himpunan !!{formula}s. Himpunan
!!{formula}s ora konsisten yen kontradiksi, kayata~$\lfalse$, iku
!!{derivable} saka himpunan kasebut; yen ora, himpunan iku konsisten.
Yen himpunan ora konsisten relatif marang sawijining sistem,
kabeh !!{formula}s ing kono ora bisa bener bebarengan ing sawijining
jagad ing model saka sistem kasebut. Kanggo teorema kasampurnan, kita
bakal mbuktekake kosok baline: saben himpunan kang konsisten relatif
marang sistem mau nduweni kabeh anggota kang bener ing sawijining jagad
ing model saka sistem kasebut, yaiku
``model kanonik.''

\begin{defn}
  Himpunan $\Gamma$ iku \emph{konsisten} relatif marang sistem~$\Sigma$
  utawa, kaya kang bakal kita sebut, $\Sigma$-konsisten, yen lan mung yen $\Gamma
  \Proves/[\Sigma] \lfalse$.
\end{defn}

Dadi, upamane, himpunan $\{ \Box(p \lif q), \Box p, \lnot\Box q \}$
konsisten relatif marang logika proposisional, nanging ora
\Log{K}-konsisten. Kanthi cara padha, himpunan $\{ \Diamond p, \Box\Diamond
p \lif q, \lnot q \}$ ora \Log{K5}-konsisten.

\begin{prop}\ollabel{prop:consistencyfacts}
  Upamane $\Gamma$ himpunan !!{formula}s. Banjur:
  \begin{enumerate}
  \item $\Gamma$ iku $\Sigma$-konsisten yen lan mung yen ana
    !!{formula}~$!A$ saengga $\Gamma \Proves/[\Sigma]
    !A$.
  \item \ollabel{prop:consistencyfacts-b}%
    $\Gamma \Proves[\Sigma] !A$ yen lan mung yen $\Gamma \cup \{
    \lnot!A \}$ ora $\Sigma$-konsisten.
  \item \ollabel{prop:consistencyfacts-c}%
    Yen $\Gamma$ $\Sigma$-konsisten, mula kanggo saben !!{formula}
    $!A$, $\Gamma \cup \{ !A \}$ iku
    $\Sigma$-konsisten utawa $\Gamma \cup \{ \lnot!A \}$ iku
    $\Sigma$-konsisten.
  \end{enumerate}
\end{prop}

\begin{proof}
  Sipat-sipat iki gampang diturunake nganggo logika proposisional klasik.
  Kita menehi argumen kanggo \olref{prop:consistencyfacts-c}. Buktekna
  kanthi kontraposisi: upamane $\Gamma \cup \{ !A \}$ lan
  $\Gamma \cup \{ \lnot!A \}$ loro-lorone ora $\Sigma$-konsisten. Banjur
  miturut definisi konsistensi, $\Gamma, !A \Proves[\Sigma]
  \lfalse$ lan $\Gamma, \lnot !A \Proves[\Sigma] \lfalse$. Miturut
  teorema deduksi, $\Gamma \Proves[\Sigma] !A \to \lfalse$ lan
  $\Gamma \Proves[\Sigma] \lnot!A \lif \lfalse$. Nanging $(!A \lif
  \lfalse) \lif ((\lnot!A \lif \lfalse) \lif \lfalse)$ iku
  instansi tautologis, mula miturut
  \olref[prp]{prop:derivabilityfacts}\olref[prp]{prop:derivabilityfacts-ruleT},
  $\Gamma \Proves[\Sigma] \lfalse$.
\end{proof}

\end{document}
